\documentclass[11pt]{article}

\newcommand{\cnum}{Math 171}
\newcommand{\ced}{Winter 2026}
\newcommand{\ctitle}[4]{\title{\vspace{-0.5in}\cnum, \ced\\Assignment #1: #2}\author{\vspace{-0.35in}\\#3, #4}}
\usepackage{enumitem}
\usepackage[usenames,dvipsnames,svgnames,table,hyperref]{xcolor}
\usepackage{amsmath, amsfonts}
\usepackage{graphicx} % Required for inserting images
\usepackage{float}
\usepackage{listings}
\usepackage{xcolor}
\usepackage{hyperref}
\usepackage{comment}

\renewcommand*{\theenumi}{\alph{enumi}}
\renewcommand*\labelenumi{(\theenumi)}
\renewcommand*{\theenumii}{\roman{enumii}}
\renewcommand*\labelenumii{\theenumii.}

\definecolor{codegreen}{rgb}{0,0.6,0}
\definecolor{codegray}{rgb}{0.5,0.5,0.5}
\definecolor{codepurple}{rgb}{0.58,0,0.82}
\definecolor{backcolour}{rgb}{0.95,0.95,0.92}
\definecolor{header}{HTML}{205459}
\definecolor{solution}{HTML}{204d52}

\newcommand{\solution}[1]{{{\textcolor{header}{Solution:} \textcolor{solution}{#1}}}}
\setlist[enumerate]{label=(\alph*)}

\lstdefinestyle{mystyle}{
    backgroundcolor=\color{backcolour},
    commentstyle=\color{codegreen},
    keywordstyle=\color{magenta},
    numberstyle=\tiny\color{codegray},
    stringstyle=\color{codepurple},
    basicstyle=\ttfamily\footnotesize,
    breakatwhitespace=false,
    breaklines=true,
    captionpos=b,
    keepspaces=true,
    numbers=left,
    numbersep=5pt,
    showspaces=false,
    showstringspaces=false,
    showtabs=false,
    tabsize=2
}


\begin{document}
\ctitle{\#1}{}{Asher Christian}{006-150-286}
\date{}
\maketitle

\section{Exercise 1.1}
\solution{
    Consider the probability
    \[
        \mathbb{P}(X_{n+1} = i | X_{n} = j, X_{n-1} = i_{n-1}, \dots, X_0 = i_0) = p(i,j)
    \] 
    However consider
    \[
    \mathbb{P}(X_3 = 0 | X_2 = 1, X_1 = 2) = \frac{1}{2}
    \] 
    since if $X_2 = 1, X_1 = 2$ then $Y_0 = 1, Y_1 = 1, Y_2 = 0$ so $Y_3$ can be either $1$ or  $0$ with equal probability
    and so  $X_3$ can be either $1$ or  $0$ with equal probability. However,
    \[
    \mathbb{P}(X_3 = 0 | X_2 = 1, X_1 = 0) = 0
    \] 
    since in this scenario we must have $Y_0 = 0, Y_1 = 0, Y_2 = 1$ with $Y_3$ being either $1$ or $0$ with equal probability thus
    $X_3 = 1 $ or $2$ with equal probability. Therefore  $X_n$ does not depend only on  $X_{n-1}$ and the sequence is not a markov chain
}

\section{Exercise 1.2}
\solution{
    We need to compute
    \[
        p(i,j) = \mathbb{P}(X_{n+1} = j | X_n = i)
    \] 
    if $i = 5$ then  $j = 4$ with probability 1. Similalry if  $i=0$ then  $j=1$ with probability  $1$.
    if $i \in [1,4]$ then  $ \mathbb{P}(X_{n+1} =i-1) = \frac{i}{5}\frac{i}{5}$ and $ \mathbb{P}(X_{n+1} = i+1) = \frac{5-i}{5}(\frac{5-i}{5})$ with all the other probability being that $X_{n+1} = X_n$
    Then
    \[
    P :=
    \begin{pmatrix} 0 & 1 & 0 & 0 & 0 & 0\\
    \frac{1}{25} & \frac{8}{25} & \frac{16}{25} & 0 & 0 & 0 \\
    0 & \frac{4}{25} & \frac{12}{25} & \frac{9}{25} & 0 & 0\\
    0 & 0 & \frac{9}{25} & \frac{12}{25} & \frac{4}{25} & 0\\
    0 & 0 & 0 & \frac{16}{25} & \frac{8}{25} & \frac{1}{25}\\
    0 & 0 & 0 & 0 & 1 & 0
    \end{pmatrix} 
    \] 
    and
    \[
        p(i,j) = P_{i,j}
    \] 
}

\section{Exercise 1.3}
\solution{
    The probabilities of the rolls are
    \begin{align*}
        2 & \qquad\frac{1}{16}\\
        3 & \qquad\frac{2}{16}\\
        4 & \qquad\frac{3}{16}\\
        5 & \qquad\frac{4}{16}\\
        6 & \qquad\frac{3}{16}\\
        7 & \qquad\frac{2}{16}\\
        8 & \qquad\frac{1}{16}
    \end{align*}
    grouping by value modulo 6 we have the probabilities for the increase as follows
    \begin{align*}
        0 & \qquad\frac{3}{16}\\
        1 & \qquad\frac{2}{16}\\
        2 & \qquad\frac{2}{16}\\
        3 & \qquad\frac{2}{16}\\
        4 & \qquad\frac{3}{16}\\
        5 & \qquad\frac{4}{16}
    \end{align*}
    putting these into a matrix we get
    \[
    P := \frac{1}{16}
    \begin{pmatrix} 
        3 & 2 & 2 & 2 & 3 & 4\\
        4 & 3 & 2 & 2 & 2 & 3\\
        3 & 4 & 3 & 2 & 2 & 2\\
        2 & 3 & 4 & 3 & 2 & 2\\
        2 & 2 & 3 & 4 & 3 & 2\\
        2 & 2 & 2 & 3 & 4 & 3
    \end{pmatrix} 
    \] 
}

\section{Exercise 1.4}
for each human $X$ let $X =1$ signify  $X$ is a homeowner and  $X = 0$ mean  $X$ is a renter and $n$ be the $n$-th decade then
\[
    \mathbb{P}(X_{n+1} = 1| X_n = 0) = \frac{12}{100}
\] 
\[
    \mathbb{P}(X_{n+1} = 0 | X_n = 0) = \frac{88}{100}
\] 
\[
    \mathbb{P}(X_{n+1} = 0 | X_n = 1) = \frac{6}{100}
\] 
\[
    \mathbb{P}(X_{n+1} = 1 | X_n = 1) = \frac{94}{100}
\] 
so the transition matrix is
\[
P :=
\begin{pmatrix} 
    \frac{88}{100} & \frac{12}{100} \\
    \frac{6}{100} & \frac{94}{100}
\end{pmatrix} 
\] 
Now $X$ has starting distribution $\begin{pmatrix} \frac{64}{100} & \frac{36}{100} \end{pmatrix} $ in 1990.
THen $XP = \begin{pmatrix} .5848 & .4152 \end{pmatrix} $ is the distribution in $2000$
and $XP^2 = \begin{pmatrix} .539536 & .460464 \end{pmatrix} $ is the distribution in $2010$ so in conclusion, 
in $2010$  $46$ percent of households were homeowners

\section{Exercise 1.5}

First lets construct the transition matrix
 \[
     P_{i,j} = p(i,j)
\] 
\[
 P =
 \begin{pmatrix} 1 & 0 & 0 & 0 & 0\\
     .6 & 0 & .4 & 0 & 0\\
     0 & .6 & 0 & .4 & 0\\
     0 & 0 & .6 & 0 & .4\\
     0 & 0 & 0 & 0 & 1
 \end{pmatrix} 
\] 
then we can compute $p^{3}(1,4)$ and $p^{3}(1,0)$ in two different ways
first we can compute $P^{3}$ 
\[
P^{3} = \begin{pmatrix} 
    1 & 0 & 0 & 0 & 0\\
    0.744 & 0 & 0.192 & 0 & 0.064\\
    0.36 & 0.277 & 0 & 0.192 & 0.16\\
    0.216 & 0 & 0.288 & 0 & 0.496\\
    0 & 0 & 0 & 0 & 1
\end{pmatrix} 
\] 
reading directly we get
\[
p^{3}(1,4) = 0.064
\] 
and
\[
p^{3}(1,0) = 0.744
\] 
we can also compute directly
\[
    p^{3}(1,4) = \mathbb{P}(X_4 = 4 | X_1 = 1) = \sum_{}^{} \mathbb{P}(X_2 = k | X_1 =1) \mathbb{P}(X_3 = j | X_2 = k) \mathbb{P}(X_4 = 4 | X_3 = j)
\] 
over all possible combinations of $k,j$
\[
= \mathbb{P}(X_2=2|X_1=1) \mathbb{P}(X_3=3|X_2=2) \mathbb{P}(X_4=4 |X_3=3) = p(1,2)p(2,3)p(3,4) = .4^{3} = .064
\] 
additionally we can compute $p^{3}(1,0)$ by computing all possible paths from $1$ to $0$ in three steps. The first is that  $X_2 = 0$ and it stays there,
the second is $1-2-1-0$ and those are the only paths their probabilities are computed and summed
 \[
 \mathbb{P}(X_2 = 0 | X_1 = 1) + \mathbb{P}(X_2 = 2| X_1=1) \mathbb{P}(X_3 = 1| X_2 = 2) \mathbb{P}(X_4 = 0 | X_3 = 1) 
\]\[
=p(1,0) + p(1,2)p(2,1)p(1,0) = .6 + .4 * .6 * .6 = .744
\] 
so our answers coincide.

\section{Exercise 1.6}
Let state A be 1, state B be 2, and state C be 3 then we construct the transition matrix
\[
P =
\begin{pmatrix} 0 & \frac{1}{2} & \frac{1}{2} \\
    \frac{3}{4} & 0 & \frac{1}{4}\\
    \frac{3}{4} & \frac{1}{4} & 0
\end{pmatrix} 
\] 
Lets compute the higher powers to answer the rest of the question
\[
    P^{2} = \begin{pmatrix}
.75 & .125 & .125\\
.1875 & .4375 & .375\\
0.1875 & 0.375 & 0.4375
\end{pmatrix}
\qquad P^{3} =
\begin{pmatrix} 0.1875 & 0.40625 & 0.40625\\
    0.609375 & 0.1875 & 0.203125\\
    0.609375 & 0.203125 & 0.1875
\end{pmatrix} 
\] 
assuming the taxicab driver begins at the airport at time 0 the probability he is at the airport at time 2 
is $0.75$ and he is at the other hotels with equal probability  $0.125$. The probability he is at
hotel  $B$ or state 2 at time $3$ is $0.40625$

\section{Exercise 1.7}
I map the states as follows  $(RR,RS,SR,SS) \rightarrow (1,2,3,4)$ and create the transition matrix
\[
P =
\begin{pmatrix} 0.6 & 0.4 & 0 & 0\\
    0 & 0 & 0.6 & 0.4\\
    0.6 & 0.4 & 0 & 0\\
    0 & 0 & 0.3 & 0.7
\end{pmatrix} 
\qquad P^2 =
\begin{pmatrix} 
    0.36 & 0.24 & 0.24 & 0.16\\
    0.36 & 0.24 & 0.12 & 0.28\\
    0.36 & 0.24 & 0.24 & 0.16\\
    0.18 & 0.12 & 0.21 & 0.49
\end{pmatrix} 
\] 
if $X_1 = (SS)$ where day 1 is monday, then we can compute $X_3 = (0,0,0,1) P^2 = (0.18,0.12,0.21,0.49)$ so the probability that it rains on
Wednesday is $0.18 + 0.21 = 0.39$

\end{document}
